\documentclass[10pt,a4paper]{amsart}
\usepackage[margin=19mm]{geometry}
\usepackage{amsmath,amssymb,mathtools}
\usepackage[scr=rsfs,cal=euler]{mathalfa}
\usepackage{xfrac}
\usepackage[unicode,hidelinks,bookmarks=false]{hyperref}
\newcommand{\ie}{i.e.\ }
\newcommand{\cf}{cf.\ }
\newcommand{\resp}{respectively\ }
\newcommand{\Subsection}{\subsection}
\providecommand{\nobreakdash}{\nobreak\hbox{-}}
\setlength{\emergencystretch}{2em}
\multlinegap0pt
\usepackage{amssymb}
\newtheorem{proposition}{Proposition}
\title{A remark about Calder\'on-Hardy spaces with variable exponents}
\author{Pablo Rocha}
\date{Source: arXiv:2512.17175v1 (CC BY 4.0)}
\begin{document}
\maketitle

\noindent\emph{Source and licence.} A remark about Calder\'on-Hardy spaces with variable exponents. Version arXiv:2512.17175v1 (CC BY 4.0), distributed by arXiv under the Creative Commons Attribution 4.0 International licence, http://creativecommons.org/licenses/by/4.0/, as recorded in the arXiv licence metadata for this exact version and shown on the abstract page https://arxiv.org/abs/2512.17175v1. Full licence notice: \texttt{licenses/arxiv-2512.17175-CC-BY-4.0.txt}.\par\smallskip

\noindent\textit{Editorial scope.} Counted unit: the article's proposition functions (source0.tex lines 186--201). The article's complete original proof of it is delivered with it (source0.tex lines 203--205, one \texttt{proof} environment of 3 lines). No mathematics is written, repaired or re-derived: the statement and the proof are the article's own lines, in the article's own notation, and no retained line is substituted or re-flowed.  No further source-local context is needed: every cross-reference of every retained line resolves inside the two delivered spans.  Nothing was omitted from any retained span, so the package carries no omission record.  The retained lines cite 1 works, whose entries are transcribed verbatim from the article's own bibliography source in the e-print and delivered with short editorial labels recorded in the item's reference mapping.  No retained line contains a figure environment, an image-inclusion command or drawing code, so no image is delivered and the package declares no asset dependency.  Editorial formatting only, in addition: an AMS \texttt{amsart} preamble that restates the article's own declarations and macros used by the retained lines, namely the environments \texttt{proposition} on separate counters, together with the standard packages those lines need.  The retained environments print in this document as Proposition 1 in order of appearance, whereas the article prints the same items with the numbering of its own full document.  The complete native source of the article as distributed in the arXiv e-print for this exact version is bundled unchanged as \texttt{sources/191298/source0.tex}.
\begin{proposition} \label{b_k functions}
Let $p(\cdot) : \mathbb{R}^{n} \to (0, \infty)$ such that $p(\cdot) \in LH_{0} \cap LH_{\infty}(\mathbb{R}^{n})$ and 
$0 < p_{-} \leq p_{+} < \infty$. Let $s > 1$ and $0 < p_{*} < \underline{p}$ such that $s p_{*} > p_{+}$ and let 
$\{ b_j \}_{j=1}^{\infty}$ be a sequence of nonnegative functions in $L^{s}(\mathbb{R}^{n})$ such that each $b_j$ is supported 
in a cube $Q_j \subset \mathbb{R}^{n}$ and
\begin{equation} \label{bks}
\| b_j \|_{L^{s}(\mathbb{R}^{n})} \leq A_j |Q_j|^{1/s},
\end{equation}
where $A_j >0$ for all $j \geq 1$. Then, for any sequence of nonnegative numbers $\{ k_j \}_{j=1}^{\infty}$ we have
\[
\left\| \sum_{j=1}^{\infty} k_j b_j \right\|_{L^{p(\cdot)/p_{*}}(\mathbb{R}^{n})} \leq C \left\| \sum_{j=1}^{\infty} A_j k_j 
\chi_{Q_j} \right\|_{L^{p(\cdot)/p_{*}}(\mathbb{R}^{n})},
\]
where $C$ is a positive constant which does not depend on $\{ b_j \}_{j=1}^{\infty}$, $\{ A_j \}_{j=1}^{\infty}$, and 
$\{ k_j \}_{j=1}^{\infty}$.
\end{proposition}

\begin{proof}
The proof is similar to the one given in \cite[Proposition 3.3]{Rocha2}.
\end{proof}

\begin{thebibliography}{99}
\bibitem[Rocha2]{Rocha2} {\sc P. Rocha}, \textit{Convolution operators and variable Hardy spaces on the Heisenberg group}, Acta Math. Hung. 174 (2) (2024), 429-452.
\end{thebibliography}
\end{document}
